\section{Filtered direct images on simple normal-crossing supports}\label{sec:hodge}

The residue insertion has placed the finite lifting obstruction in
\(R^1h_*\omega_H\).  We need two facts about this absolute dualizing
direct image: it embeds as the lowest step of a filtered right
\(\Dcal\)-module, and on a projective base the kernel of first-symbol
multiplication admits no map from an ample line.  In
Section~\ref{sec:lifting}, differentiation of representatives of the
finite obstruction will produce just such a symbol relation.  We prove
the two facts here at the geometric scope of the graph supports.

Let \(p:A\to T\) be a holomorphic map of complex manifolds, where \(A\)
has pure dimension \(n\).  Let \(i:H\hookrightarrow A\) be a reduced
closed analytic subspace of pure dimension \(d\), and put \(c=n-d>0\).
Assume the following.
\begin{enumerate}
 \item The map \(p\) is a submersion near \(H\), and \(h=p|_H:H\to T\)
 is proper.
 \item Locally in \(A\), the space \(H\) is a reduced SNC divisor in a
 smooth submanifold of dimension \(d+1\).  It has finitely many globally
 smooth indexed irreducible components \(H_1,\ldots,H_s\), and every
 closed intersection \(H_I=\bigcap_{i\in I}H_i\) is smooth, possibly
 empty or disconnected.
 \item Every connected component of every nonempty \(H_I\) is proper
 over \(T\) and has a global closed real \((1,1)\)-form which, locally
 on \(T\), becomes K\"ahler after adding the pullback of a K\"ahler form
 from \(T\).
\end{enumerate}
The last condition is the relative-form convention for a proper K\"ahler
map.  A K\"ahler form on a neighborhood of \(H\), as constructed in the
preceding section, implies it.  We may replace \(A\) by the submersion
neighborhood of \(H\).  All direct images below are proper on their
support; the ambient map \(p\) itself need not be proper.

We use right \(\Dcal\)-modules with the increasing order filtration.
Set
\begin{equation}\label{eq:support-module}
 K=\mathcal H^c_{[H]}(\omega_A),\qquad
 F_0K=\im\bigl(\mathcal{E}xt^c_A(\OO_H,\omega_A)\longrightarrow K\bigr),
 \qquad F_pK=(F_0K)F_p\Dcal_A\quad(p\geq0),
\end{equation}
and \(F_pK=0\) for \(p<0\).  The brackets mean algebraic local
cohomology with finite pole orders.  We will show that the first map is
injective and its image is the absolute dualizing sheaf \(\omega_H\).

For filtered right \(\Dcal\)-modules, \(p_+\) denotes the direct image;
in this submersion setting it is computed by the relative right de Rham
(Spencer) complex followed by \(Rp_*\).

\begin{proposition}[The filtered SNC direct image]\label{prop:filtered-snc}
For the preceding data, \((K,F)\) is a good filtered regular holonomic
right module.  Put \(M^j=\mathcal H^j p_+K\) for any integer \(j\), and
give it the filtration induced by the filtered direct image.  Every
level cohomology map
\[
 \mathcal H^j F_p p_+(K,F)\longrightarrow M^j
\]
is injective, with image \(F_pM^j\).  The module \((M^j,F)\) has a finite
filtration by \(\Dcal_T\)-submodules, strict at every \(F_p\), whose
quotients are polarizable real pure Hodge modules.  More precisely, the
quotient with support-filtration index \(\ell\) has weight \(\ell+j\).
There are canonical identities, with the first realized by that injection,
\begin{equation}\label{eq:lowest-direct-image}
 F_0M^j=R^jh_*\omega_H\lhook\joinrel\longrightarrow M^j,
 \qquad F_pM^j=0\quad(p<0).
\end{equation}
Here \(R^j=0\) for \(j<0\), and \(\omega_H\) is the absolute dualizing
sheaf.
\end{proposition}

For the later use with \(j=1\), define the first-symbol map for the right
action by
\[
 \sigma:F_0M^1\otimes T_T\longrightarrow\gr^F_1M^1,
 \qquad m\otimes\xi\longmapsto[m\xi].
\]
When \(T\) is smooth projective, Corollary~\ref{cor:symbol-kernel} will
prove \(\operatorname{Hom}_{\OO_T}(N,\ker\sigma)=0\) for every ample line
\(N\) on \(T\).

\subsection{Finite poles and the two filtrations}

At a point of \(H\), choose analytic coordinates
\[
 (z_1,\ldots,z_{c-1},x_1,\ldots,x_t,y),\qquad
 \mathscr I_H=(z_1,\ldots,z_{c-1},f),\quad f=x_1\cdots x_t.
\]
The \(z\)-list is empty when \(c=1\).  Write \(R\) for the local analytic
ring and \(\eta\) for its coordinate volume form.  These equations are
a regular sequence.  Their localization \v Cech complex, equivalently
the direct limit of their Koszul complexes, has cohomology only in degree
\(c\).  In that degree it is
\begin{equation}\label{eq:polar-cech}
 \frac{R[z_1^{-1},\ldots,z_{c-1}^{-1},f^{-1}]\eta}
 {\displaystyle\sum_{j=1}^{c-1}
 R[z_1^{-1},\ldots,\widehat{z_j^{-1}},\ldots,z_{c-1}^{-1},f^{-1}]\eta
 +R[z_1^{-1},\ldots,z_{c-1}^{-1}]\eta}.
\end{equation}
The same assertion holds for a subunion of the branches, with the product
of its selected \(x_i\)'s.  An intersection of \(k\) branches instead
has \(c+k-1\) regular equations.

Successive finite Taylor divisions give a unique additive polar normal
form in \eqref{eq:polar-cech}: a finite sum of terms
\begin{equation}\label{eq:polar-normal-form}
 a_{\mathbf u,\mathbf v,J}(x_{J^c},y)
 \prod_{j=1}^{c-1}z_j^{-u_j}\prod_{i\in J}x_i^{-v_i}\eta,
 \qquad u_j,v_i\geq1,\quad \varnothing\ne J\subset\{1,\ldots,t\}.
\end{equation}
The coefficient is a holomorphic germ in exactly the displayed
variables.  Each fraction has bounded negative orders, so only finitely
many Taylor coefficients in its denominator variables enter this
description.  It is an additive normal form, not an assertion that the
polar types form an \(R\)-linear direct sum.

The Koszul generator of \(\mathcal{E}xt^c_A(\OO_H,\omega_A)\) maps to
the class
\begin{equation}\label{eq:simple-fraction}
 \frac{a\eta}{z_1\cdots z_{c-1}f}.
\end{equation}
If it is zero in \eqref{eq:polar-cech}, the expansion of \(a|_{z=0}/f\)
has no negative \(x_i\)-power.  Every Taylor monomial of \(a|_{z=0}\)
is then divisible by every \(x_i\), hence by \(f\).  Thus
\(a\in(z_1,\ldots,z_{c-1},f)\), proving injectivity.  Complete-intersection
adjunction, including its determinant of the conormal, identifies its
image intrinsically with \(\omega_H\).

For a polar term define the excess order
\[
 e(\mathbf u,\mathbf v,J)=
 \sum_{j=1}^{c-1}(u_j-1)+\sum_{i\in J}(v_i-1).
\]
The simple fractions \eqref{eq:simple-fraction} are exactly the terms
of excess zero.  On a right canonical module a coordinate derivative
acts by minus differentiation of the coefficient of \(\eta\).  A
derivative in a denominator variable raises its pole order by one with
a nonzero scalar.  Conversely, in \eqref{eq:polar-normal-form} the
coefficient is independent of those variables, so these derivatives
produce each desired term from an excess-zero term without a tangential
error.  It follows that
\begin{equation}\label{eq:pole-filtration}
 F_pK=\{\text{finite sums of terms \eqref{eq:polar-normal-form} of excess
 at most }p\},\qquad F_pK=0\ (p<0).
\end{equation}
This proves the generated formula in \eqref{eq:support-module} locally,
and proves that it is a good filtration.  The generated definition
makes it independent of the coordinates.

For a closed immersion of smooth manifolds \(j:Y\hookrightarrow A\),
the right transfer convention is
\begin{equation}\label{eq:right-transfer}
 F_pj_+(N,F)=\sum_{\nu}F_{p-|\nu|}N\,\partial_{\perp}^{\nu}.
\end{equation}
There is no codimension shift.  For example, the codimension shift for
left transfer cancels the two dimension shifts in changing sides.  These
right conventions are recorded in
\cite[Section~1.1, equations~(1.1.2)--(1.1.4)]{FujinoFujisawaSaito2014}.

We define an increasing support filtration on \(K\) by summing the
images of the local-cohomology modules of subunions of at most \(k\)
of the indexed components:
\begin{equation}\label{eq:support-weight}
 W_{-d+k-1}K=
 \sum_{|I|\leq k}\im\left(
 \mathcal H^c_{[\bigcup_{i\in I}H_i]}(\omega_A)\longrightarrow K\right)
 \quad(1\leq k\leq s).
\end{equation}
Set the lower steps to zero and the upper steps to \(K\).  In the polar
normal form the summand for \(I\) consists of the types \(J\subset I\).
The support maps are locally injective, and hence
\[
 W_{-d+k-1}K=\{\text{terms with }|J|\leq k\}.
\]
The intersection \(F_pK\cap W_{-d+k-1}K\) imposes this bound and the
excess bound simultaneously.  Keeping exactly the polar type \(J=I\),
\(|I|=k\), leaves \(c+k-1\) normal denominators.  By
\eqref{eq:right-transfer} the resulting filtered quotient is
\begin{equation}\label{eq:weight-gradeds}
 \gr^W_{-d+k-1}(K,F)
 \simeq\bigoplus_{|I|=k}(i_I)_+(\omega_{H_I},F^{(0)}),
 \qquad
 F^{(0)}_p\omega_{H_I}=
 \begin{cases}0,&p<0,\\ \omega_{H_I},&p\geq0.
 \end{cases}
\end{equation}
Here \(i_I:H_I\hookrightarrow A\), empty intersections contribute zero,
and \(d_I=\dim H_I=d-k+1\).  Distinct sets \(I\) give distinct polar
types even at a point incident to further components.  This proves the
exact induced \(F\) in \eqref{eq:weight-gradeds}, not just the unfiltered
quotient.

The identifications glue intrinsically.  For an ordered \(I\) of size
\(k\), the iterated connecting maps for Mayer--Vietoris triangles of
supports give
\begin{equation}\label{eq:support-mayer-vietoris}
 \frac{\mathcal H^c_{[\bigcup_{i\in I}H_i]}(\omega_A)}
 {\displaystyle\sum_{J\subsetneq I}\im
   \mathcal H^c_{[\bigcup_{j\in J}H_j]}(\omega_A)}
 \longrightarrow \mathcal H^{c+k-1}_{[H_I]}(\omega_A).
\end{equation}
Locally this map selects the term with every \(i\in I\) in the polar
type, up to the sign of the order, and is therefore an isomorphism.
The one fixed ordering of the global smooth components fixes those signs
on overlaps.  Disconnected intersections are treated componentwise.
This identifies \eqref{eq:support-mayer-vietoris} with
\eqref{eq:weight-gradeds} globally.  The smooth-support modules on the
right are regular holonomic; their finite extension \(K\) is regular
holonomic as well.

\subsection{The real structure and strict direct image}

The filtration \(W\) has a real realization.  With the absolute right
de Rham complex ending in degree zero, \(\DR_A\omega_A\simeq\C_A[n]\).
The ordinary unshifted finite-pole de Rham complex along one coordinate
has the degree-zero constant and the degree-one logarithmic class of a punctured disk.
Tensoring these local comparisons and then taking the localization
\v Cech complexes gives, naturally for all inclusions of subunions,
\begin{equation}\label{eq:real-support}
 \DR_A K\simeq R\Gamma_H\C_A[n+c]
 \simeq i_*\mathbb D_H(\C_H[d]).
\end{equation}
Here \(R\Gamma_H\) is the sheaf-valued support functor on \(A\).
The shift in the second identity is fixed by the complex orientation:
\(n+c=2n-d\).  The same support complexes with \(\R\) supply a real
perverse sheaf, since complexification is exact and faithful and its
complexification is the de Rham realization of the regular holonomic
module.  Taking real perverse images of the subunion maps defines a real
filtration whose complexification is \eqref{eq:support-weight}.

The real version of \eqref{eq:support-mayer-vietoris} is an isomorphism
because its complexification is.  Apply the support comparison anew to
its graded target \(\mathcal H^{c+k-1}_{[H_I]}(\omega_A)\), now with
support dimension \(d_I\).  Its real realization is
\[
 (i_I)_*\mathbb D_{H_I}(\R_{H_I}[d_I])
 \simeq (i_I)_*\R_{H_I}[d_I],
\]
using smoothness and the complex orientation of \(H_I\).  This is the
shifted rank-one constant real local system of the graded target.
The residue normalization may multiply its complex realization
by a nonzero constant on a connected component; it introduces neither
a varying factor nor monodromy.  Such a real rank-one form is polarizable
for its single Hodge type.  With the filtration in
\eqref{eq:weight-gradeds}, the right module on \(H_I\) is the dual
constant
\[
 \R^H_{H_I}[d_I](d_I).
\]
The untwisted constant has weight \(d_I\) and first right step
\(F_{-d_I}\); the twist by \(d_I\) moves that step to zero and changes
the weight by \(-2d_I\).  Thus its weight is
\(-d_I=-d+k-1\).  These dual-constant conventions agree with
\cite[equations~(1.2.1)--(1.2.6) and Theorem~2.3]{FujinoFujisawaSaito2014}.

There is also a useful support check.  If a holomorphic germ \(g\)
vanishes on reduced \(H\), then
\begin{equation}\label{eq:filtered-support-condition}
 gF_pK\subset F_{p-1}K.
\end{equation}
Indeed \(g\in(z_1,\ldots,z_{c-1},f)\).  Multiplication by \(z_j\)
lowers the excess or kills the polar class; multiplication by \(f\)
does the same, since a surviving class has some remaining \(x_i\)-pole.
This is the filtered support condition for treating the object on \(H\)
through its local embeddings \cite[Section~1.17]{Saito1990Kahler}.
In particular properness on \(H\), rather than properness of \(A\),
is the relevant properness here.

We now prove Proposition~\ref{prop:filtered-snc}.  Apply the proper
K\"ahler direct-image theorem
\cite[Theorem~1 and Remark~1.2]{Saito2022Kahler} after restricting to
each connected component of \(T\), and then separately to the constant
real Hodge module on each connected smooth component of \(H_I\) over it,
using the stipulated global relative form.  These componentwise sums are
locally finite on \(T\): near any parameter, apply
Lemma~\ref{lem:finite-components} to the proper support and its smooth
closed intersections.  Thus the componentwise direct images and their
polarizations give those for the full \(H_I\); in the cover and graph
applications the final indexing is already finite.  The theorem supplies
strict filtered direct image, Lefschetz decomposition, and real primitive
polarizations.
The latter give a polarization on each full cohomology module.  After
twisting by \(d_I\) and using filtered closed-embedding transitivity,
the degree-\(v\) direct image of a summand of
\(\gr^W_\ell K\) is polarizable real pure of weight
\begin{equation}\label{eq:direct-image-weight}
 d_I+v-2d_I=-d_I+v=\ell+v.
\end{equation}
We use the convention fixed by the constant input in Section~1.1 and
the direct image in Section~2.2 of that source.  The opposite sign for
the untwisted dimension in its displayed theorem statement is a
typographical error, also ruled out by applying it to the identity map.
This application is only to constant modules on smooth sources.

Write \(C=p_+(K,F)\) for the filtered direct-image complex.  Since
\(p\) is a submersion near its support, before applying \(Rp_*\) its
right relative de Rham complex has at level \(F_p\) and degree \(-v\)
the term
\begin{equation}\label{eq:relative-right-dr}
 F_{p-v}K\otimes_{\OO_A}\bigwedge^v T_{A/T}.
\end{equation}
The tangent wedges are locally free.  The simultaneous pole and branch
bounds show that the quotient of this level complex by \(W\) is exactly
the same level of the relative complex for \(\gr^W K\).

Use the increasing \(W\) spectral sequence, written as
\begin{equation}\label{eq:weight-spectral-sequence}
 E_1^{-\ell,v+\ell}=\mathcal H^v p_+\gr^W_\ell K
 \quad\Longrightarrow\quad \mathcal H^v p_+K.
\end{equation}
Each term on the first page is a strict filtered direct image and a
polarizable real pure module of weight \(\ell+v\), by
\eqref{eq:direct-image-weight}.  Strictness says that the corresponding
level-\(F_p\) first page injects into it with image \(F_p\).
The real support realization \eqref{eq:real-support}, its proper direct
image on the support, and the de Rham comparison make the unfiltered
differentials real.  Comparison with the level spectral sequences makes
them filtration preserving.

The differential \(d_r\) sends \((\ell,v)\) to
\((\ell-r,v+1)\).  For \(r=1\) these weights are equal.  The differential
is therefore a morphism of real pure Hodge modules, hence is strict for
\(F\), and its kernels and cokernels remain polarizable pure modules
\cite[Propositions~1.10 and~1.14]{Saito1990Kahler}.  Consequently the
level second page injects into the unfiltered second page with image
\(F_p\), and the latter consists of pure pieces of the indicated weights.

For \(r\geq2\) the target weight \(\ell+v-r+1\) is smaller than the
source weight \(\ell+v\).  Here is why a real filtered differential
between those pieces is zero.  Decompose both into their locally finite
strict-support summands.  A perverse map between distinct strict supports
has image supported properly in at least one of them, and hence is zero.
On a common smooth dense support of dimension \(e\), the map is a real
map of local systems preserving their Hodge filtrations; this follows
also directly from \eqref{eq:right-transfer}, which recovers the intrinsic
filtered module as the sections annihilated by the normal ideal.  The
two variation weights are the module weights minus \(e\).  A real map
preserving \(F^\bullet\) preserves \(\overline F^\bullet\) as well.
The image of a vector of source type \((a,b)\) lies in
\(F^a\cap\overline F^b\) of the target.  Since \(a+b\) is the larger
source weight, that intersection in the pure target is zero.  Thus the
generic map is zero, and strict support makes the map zero everywhere.

Inductively, injection of a level page implies its differential is zero
when the unfiltered differential is zero.  Hence both sequences
degenerate at the second page and the level infinity page injects with
image \(F_p\).  The filtration \(W\) is finite.  If the abutment map
\(\mathcal H^j(F_pC)\to\mathcal H^j(C)\) had a nonzero kernel element,
choose the least \(\ell\) containing it.  Its nonzero leading class in
\(\gr^W_\ell\) would contradict the injection on the infinity page.
This proves the asserted injection of every level.

The same leading-class argument proves strictness of the induced \(W\):
if an element of \(\mathcal H^j(F_pC)\) maps into \(W_\ell M^j\), it
already lies in \(W_\ell\mathcal H^j(F_pC)\).  The resulting finite
filtration of \((M^j,F)\) is therefore strict for every level and has
exactly the pure second-page quotients, of weights \(\ell+j\).
It follows as well that \(M^j\) is regular holonomic and its filtration
is good: these properties are preserved under this finite filtered
extension.  Taking \(\gr^F\), and then \(\gr^F\DR_T\), preserves its
short exact sequences.  If indexed by their pure weights, the output
weight filtration is the shifted filtration \(W[j]\).

For \(p<0\), every term in \eqref{eq:relative-right-dr} vanishes.
For \(p=0\), its only term is \(F_0K=i_*\omega_H\) in degree zero.
The just-proved level injection consequently gives exactly
\eqref{eq:lowest-direct-image}.  This completes the proof of
Proposition~\ref{prop:filtered-snc}.  A relative dualizing sheaf would
appear only after tensoring by \(\omega_T^{-1}\); the line in
\eqref{eq:lowest-direct-image} is absolute.

The first residue cohomology \(R^1h_*\omega_H\) is therefore embedded
as \(F_0M^1\).  For lifting, it remains to prove the symbol-kernel
vanishing when \(T\) is smooth projective.

\subsection{Projective vanishing for the pure quotients}

The projective vanishing we will use is stated for the filtered components
of the real or complex theory of Sabbah--Schnell.  The K\"ahler
direct-image theorem above supplied real polarizations; a rational
polarization was not part of its output.  We therefore need vanishing
in the real theory and must identify the actual filtration, not merely
the underlying regular holonomic module.  The following comparison does
this for each strict-support pure summand.

\begin{lemma}[Pure-component comparison and vanishing]\label{lem:real-pure-vanishing}
Let \(T\) be smooth projective, and let \((Q,F)\) be the right filtered
module of one strict-support summand of a polarizable real pure quotient
in Proposition~\ref{prop:filtered-snc}.  Write \(Z\) for its support and
\(w\) for its weight.  It is isomorphic, as a right filtered module, to
the prime holomorphic component of a pure object in
\(\mathrm{pHM}_Z(T,\R,w)\) of Sabbah--Schnell, attached to the same
generic polarized real variation.  In particular, for every ample line
\(N\) on \(T\), every integer \(p\), and every \(v<0\),
\begin{equation}\label{eq:real-negative-vanishing}
 \mathbb H^v\bigl(T,N^{-1}\otimes\gr^F_p\DR_T Q\bigr)=0.
\end{equation}
\end{lemma}
\begin{proof}
We use Saito's analytic category of pure real Hodge modules
\cite[Sections~1.4 and~1.8]{Saito1990Kahler}, which permits discrete real
indices for the \(V\)-filtration.  The pure properties supplied above
are in that category.  If the constant-source theorem is read with the
stronger rationally indexed pure conditions, those imply the broad real
specializability conditions, and induction on support dimension gives
the pure clauses of Definition~1.8 as well.
On a smooth dense analytic Zariski open \(Z^\circ\) of \(Z\), the piece
\(Q\) gives a polarized real variation \(\mathbb H\) of weight
\(w-\dim Z\), by Lemmas~1.9 and~1.13 of that source.
Theorem~16.2.1 and Corollary~16.3.6 of
\cite{SabbahSchnell2026} give a pure object
\(B\in\mathrm{pHM}_Z(T,\R,w)\) extending this same variation with its
real polarization.  Use its prime holomorphic filtered component.

Both underlying complex regular holonomic modules are the intermediate
extension \(\operatorname{IC}_Z(\mathbb H_\C)\).  For the Saito filtered
module use strict support and regularity in Remark~1.11 of
\cite{Saito1990Kahler}; for the Sabbah--Schnell component use
\cite[Section~14.2.13 and Theorem~14.7.1]{SabbahSchnell2026}.
Regular Riemann--Hilbert and uniqueness of intermediate extension
\cite[Theorem~13.2.11 and Section~13.2.12]{SabbahSchnell2026}
give the unique unfiltered identity extending the chosen identity on
\(Z^\circ\).  It remains to prove that this identity preserves \(F\)
at the boundary.

On a smooth support \(U\) of dimension \(e\), write \(\mathscr H\) for
the flat holomorphic bundle of the variation, with decreasing Hodge
filtration.  Its right filtration is
\begin{equation}\label{eq:right-generic-filtration}
 F_p(\omega_U\otimes\mathscr H)=\omega_U\otimes F^{-p-e}\mathscr H.
\end{equation}
Indeed the left filtration is \(F_p^L=F^{-p}\), and the side change is
\(F_p^R=\omega_U\otimes F_{p+e}^L\).  The same formula is recorded in
\cite[Section~16.3.12]{SabbahSchnell2026}.  The right closed transfer
\eqref{eq:right-transfer} introduces no further shift.  Thus the two
right filtrations agree away from the boundary of \(Z^\circ\).

Work locally on \(T\).  Choose a holomorphic \(g\) whose zero set
contains that boundary but no local component of \(Z\), and split into
the finitely many local strict-support factors if necessary.  Such a
\(g\) exists by prime avoidance: the boundary has smaller dimension
than each local support branch.  If the boundary is empty there is
nothing to prove.  Use the graph embedding in \(T\times\C_t\), and let
\(Q_g\) denote the transferred module.  Put \(b=\dim T\).  Saito's
left convention uses increasing \(F^L\) and decreasing \(V_L^\alpha\)
with \(t\partial_t-\alpha\) nilpotent on its \(\alpha\)-grade.  On the
graph ambient the conversion is
\begin{equation}\label{eq:graph-side-change}
 F_p^R=\omega_{T\times\C}\otimes F^L_{p+b+1},
 \qquad V_L^\alpha\longleftrightarrow V^R_{-\alpha-1}.
\end{equation}
The second rule follows because transposition sends
\(t\partial_t\) to \(-t\partial_t-1\).  In particular
\(V_L^{>-1}\) becomes \(V^R_{<0}\).  The sign change of a normal
\(\partial_t\) does not change a generated sum.

We verify the filtered surjectivity needed in Saito's reconstruction.
Put \(\Psi=\psi_{g,1}Q\) and \(\Phi=\phi_{g,1}Q\).  Strict support
gives a surjective \(\mathrm{can}:\Psi\to\Phi\) and an injective
\(\mathrm{Var}\), by \cite[Proposition~1.5]{Saito1990Kahler}.  Forgetting
\(F\), this is the nilpotent middle-extension monodromy quiver of
holonomic modules.  Its
monodromy filtrations centered at zero obey
\begin{equation}\label{eq:can-monodromy}
 \mathrm{can}(M_k\Psi)=M_{k-1}\Phi
\end{equation}
by \cite[Lemma~3.3.13]{SabbahSchnell2026}.  The centers of the weight
filtrations in Saito's pure definition are \(w-1\) for \(\Psi\) and
\(w\) for \(\Phi\).  Thus \eqref{eq:can-monodromy} says
\(\mathrm{can}(W_i\Psi)=W_i\Phi\).  Condition~(1.8.2) of Definition~1.8 and the
direct-factor Lemma~1.15 of \cite{Saito1990Kahler} put both
\(\gr_i^W\Psi\) and \(\gr_i^W\Phi\) in Saito's pure real category of
weight \(i\).  The induced surjection on each grade is a filtered pure
morphism, hence is \(F\)-strict by its Proposition~1.10.

At a stalk, take \(y\in F_p\Phi\) and choose \(i\) with \(y\in W_i\Phi\).
Strictness on the \(i\)th grade lifts its class by an element of
\(F_p\Psi\cap W_i\Psi\).  Subtract its image and repeat on the next
lower weight.  The finite monodromy filtration terminates this procedure.
It proves
\[
 \mathrm{can}(F_p\Psi)=F_p\Phi\quad\text{for all }p.
\]
The finite lifting argument obtains filtered surjectivity from pure
strictness on each weight grade.

Saito's filtered middle-extension formula
\cite[equation~(1.7.2)]{Saito1990Kahler} now applies: the local pure
factor satisfies its specializability conditions, the zero set of \(g\)
does not contain its support, and \(\mathrm{can}\) is filtered
surjective.  After \eqref{eq:graph-side-change}, that formula is
\begin{equation}\label{eq:common-middle-extension}
 F_pQ_g^R=
 \sum_{v\geq0}
 \left(V^R_{<0}Q_g^R\cap
 (j_{\mathrm{op}})_*F_{p-v}(Q_g^R|_{t\ne0})\right)\partial_t^v,
\end{equation}
where \(j_{\mathrm{op}}:T\times\C^*\hookrightarrow T\times\C\) and the
intersection is inside \((j_{\mathrm{op}})_*j_{\mathrm{op}}^{-1}Q_g^R\).
The strict-support module and its \(V\)-pieces embed there by restriction.

For the Sabbah--Schnell component, strict real specializability, pure
support, and strictness of \(\mathrm{can}\) follow respectively from
Definition~14.2.2, Theorem~14.2.19, and Corollary~14.2.23 of
\cite{SabbahSchnell2026}.  Its unfiltered intermediate-extension property
makes \(\mathrm{can}\) surjective, so it is a filtered middle extension.
Definition~10.6.1, Remark~10.6.2, and Proposition~10.6.5 of the same
source give precisely \eqref{eq:common-middle-extension} for its right
component: for \(\alpha<0\), the filtered \(V_\alpha\) is the
intersection with the filtration off the divisor, and the full module
is generated from \(V_{<0}\) by normal derivatives.

The unfiltered intermediate-extension identity preserves the canonical
\(V\)-filtration and, by \eqref{eq:right-generic-filtration}, preserves
\(F\) off the divisor.  The two right sides of
\eqref{eq:common-middle-extension} are therefore identical.  This proves
equality of \(F\) on the graph.  Formula \eqref{eq:right-transfer}
recovers the module before graph transfer, and its filtration, as the
sections annihilated by the normal ideal.  Hence the identity is filtered
on \(T\).  The local identities glue by uniqueness of the unfiltered
identity.  The comparison uses the actual generic variation
and its filtration; the two theories' internal twist symbols need not
have the same weight convention.

A pure object of Sabbah--Schnell is an object of \(\mathrm{WHM}(T)\)
with its one-step weight filtration, by Section~14.2.14 of
\cite{SabbahSchnell2026}.  Its Theorem~16.3.10, on the smooth projective
\(T\), gives negative-ample graded de Rham vanishing for either filtered
component.  Sections~8.4.1--8.4.9 of that source identify the perverse
right de Rham convention with the Spencer complex ending in the module
in degree zero, with its filtration grading preserved.  The theorem
therefore gives \eqref{eq:real-negative-vanishing} for every homogeneous
degree \(p\) in our convention.  This proves the lemma.
\end{proof}

\begin{corollary}[No ample line in the first symbol kernel]\label{cor:symbol-kernel}
In Proposition~\ref{prop:filtered-snc}, assume in addition that \(T\)
is smooth projective, and fix \(j\).  Put \(M=M^j\), and let
\[
 \sigma:F_0M\otimes T_T\longrightarrow\gr^F_1M,
 \qquad m\otimes\xi\longmapsto[m\xi]
\]
be principal-symbol multiplication for the right action.  For every
ample line \(N\) on \(T\),
\begin{equation}\label{eq:symbol-hom-vanishing}
 \operatorname{Hom}_{\OO_T}(N,\ker\sigma)=0.
\end{equation}
This assertion includes a coherent torsion subsheaf of the kernel.
\end{corollary}
\begin{proof}
On compact \(T\) the locally finite strict-support decomposition of
each pure quotient of Proposition~\ref{prop:filtered-snc} is finite.
Apply Lemma~\ref{lem:real-pure-vanishing} to each factor.  The finite
filtration of \(M\) is \(F\)-strict, so its short exact sequences remain
exact after \(\gr^F\DR_T\).  Their hypercohomology long exact sequences
give \eqref{eq:real-negative-vanishing} for \(M\) itself.  This deduction
uses the comparison for the pure quotients and the strict filtration of
\(M\).

By \eqref{eq:lowest-direct-image}, the \(p=1\) right graded de Rham
complex has exactly two terms:
\begin{equation}\label{eq:first-symbol-complex}
 \gr^F_1\DR_T M=
 \left[F_0M\otimes T_T\xrightarrow{\ \sigma\ }\gr^F_1M\right]
 \quad\text{in degrees }-1,0.
\end{equation}
Since \(N\) is invertible, its tensor is exact, and negative
hypercohomology gives
\[
 0=\mathbb H^{-1}\bigl(T,N^{-1}\otimes\gr^F_1\DR_T M\bigr)
 =H^0(T,N^{-1}\otimes\ker\sigma)
 =\operatorname{Hom}_{\OO_T}(N,\ker\sigma).
\]
The two-term calculation uses no local freeness of \(\ker\sigma\), so
it applies to torsion as well.
\end{proof}
